\documentclass[11pt,a4paper]{article}

\usepackage[utf8]{inputenc}
\usepackage[T1]{fontenc}
\usepackage[english]{babel}
\usepackage{geometry}
\usepackage{graphicx}
\usepackage{booktabs}
\usepackage{array}
\usepackage{float}
\usepackage{caption}
\usepackage{subcaption}
\usepackage{amsmath}
\usepackage{amssymb}
\usepackage{siunitx}
\usepackage{hyperref}
\usepackage{xcolor}
\usepackage{enumitem}
\usepackage{microtype}
\usepackage{xparse}

\geometry{margin=2.2cm}
\graphicspath{{Images/}}

\hypersetup{
  colorlinks=true,
  linkcolor=blue!50!black,
  urlcolor=blue!50!black,
  citecolor=blue!50!black
}

\sisetup{detect-all=true}

% Safe figure inclusion: if a file is missing, the document still compiles
% and displays a visible placeholder with the expected file path.
\NewDocumentCommand{\safeincludegraphics}{O{}m}{%
  \IfFileExists{#2}{%
    \includegraphics[#1]{#2}%
  }{%
    \fbox{%
      \begin{minipage}[c][0.22\textheight][c]{0.82\linewidth}
      \centering
      Missing figure file:\\[0.5em]
      \texttt{\detokenize{#2}}
      \end{minipage}%
    }%
  }%
}

\newcommand{\logit}{\operatorname{logit}}

\title{Moss Distribution in the LEO1 Bath:\\
Topographic Controls and Evidence for a Bounded Moisture Optimum}
\author{}
\date{28 April 2026}

\begin{document}

\maketitle

\begin{abstract}
This technical note analyzes the spatial distribution of moss in the LEO1 bath
using three rectified stitched images, a soil-surface digital elevation model,
and a container-geometry raster. The analysis is centered on the biological
pattern itself, rather than on the visually prominent drainage network alone.
Across the three dates, the moss footprint is largely persistent, whereas the
November-to-December change is interpreted mainly as seasonal greening of
previously occupied regions. The strongest stable moss signal is displaced away
from the deepest central trough and is concentrated on off-center shoulder zones.
This pattern is consistent with a bounded moisture optimum: moss is favored in
zones that remain moist, but not in zones that are maximally wet, strongly
flushed, or hydrologically unstable.
\end{abstract}

\section{Objective and Data}

The goal of this note is to determine how stable moss occurrence is related to
bath geometry and derived topographic variables. The central question is not
only where water accumulates, but why the strongest moss signal is not observed
in the deepest part of the central trough.

The analysis uses three stitched observations of the same LEO1 bath, acquired
on 3 November 2022, 22 November 2022, and 7 December 2022. These images were
interpreted together with the soil-surface raster \texttt{LEO1\_soil.asc} and
the container-geometry raster \texttt{LEO1\_container.asc}. The image data were
rectified to the bath coordinate system before comparison with the topographic
rasters.

The working hypothesis tested here is that moss responds to a moisture window,
not to wetness monotonically. In this interpretation, the central trough may be
wet, but it may also be too saturated, too frequently flushed, or too unstable
to support persistent moss growth.

\section{Image Registration and Moss Proxy}

The three RGB mosaics were cropped to the soil surface, corrected by a
perspective transform, and rescaled to a common rectangular domain. This step
removes the surrounding frame and makes the three dates directly comparable in
bath coordinates.

\begin{figure}[H]
  \centering
  \begin{subfigure}[t]{0.31\linewidth}
    \centering
    \safeincludegraphics[height=0.50\textheight,keepaspectratio]{Images/LEO_2022_11_03_cropped_rectified_rectangle.jpg}
    \caption{3 November 2022}
  \end{subfigure}
  \hfill
  \begin{subfigure}[t]{0.31\linewidth}
    \centering
    \safeincludegraphics[height=0.50\textheight,keepaspectratio]{Images/LEO_2022_11_22_cropped_rectified_rectangle.jpg}
    \caption{22 November 2022}
  \end{subfigure}
  \hfill
  \begin{subfigure}[t]{0.31\linewidth}
    \centering
    \safeincludegraphics[height=0.50\textheight,keepaspectratio]{Images/LEO_2022_12_07_cropped_rectified_rectangle.jpg}
    \caption{7 December 2022}
  \end{subfigure}
  \caption{Rectified source images used in the analysis. The large-scale moss
  distribution is already visible on the first date and remains spatially
  coherent through December. The later images show stronger green coloration,
  consistent with seasonal greening of an existing pattern.}
  \label{fig:rectified-images}
\end{figure}

A broad moss proxy was used instead of a narrow green threshold. This choice is
important because moss color changes from yellow-green to greener tones over
the observation period. The proxy therefore aims to represent moss presence
rather than a single color state. From the date-specific proxies, two derived
quantities were computed:
\begin{enumerate}[leftmargin=1.5em]
  \item a persistent moss footprint, defined by repeated occurrence across the
  three dates;
  \item a seasonal greening signal, defined by the change between early November
  and early December.
\end{enumerate}

The persistent footprint is the primary response variable in the analysis below.
It is compared with lateral position \(|x|\), downstream coordinate \(s\), trough
depth, a flow-accumulation proxy, wall proximity, and soil thickness relative to
the container.

\section{Observed Moss Pattern}

The persistent moss footprint is broad and structured. It is not concentrated in
the deepest central corridor. Instead, the main stable patches occur on the
flanks and in the outer half of the bath.

\begin{figure}[H]
  \centering
  \begin{subfigure}[t]{0.46\linewidth}
    \centering
    \safeincludegraphics[width=0.82\linewidth]{Images/persistent_moss_map.jpg}
    \caption{Persistent moss footprint}
  \end{subfigure}
  \hfill
  \begin{subfigure}[t]{0.46\linewidth}
    \centering
    \safeincludegraphics[width=0.82\linewidth]{Images/seasonal_greening_map.jpg}
    \caption{Seasonal greening}
  \end{subfigure}
  \caption{Moss-centered diagnostics. The stable moss footprint is displaced
  from the deepest central trough, whereas the strongest seasonal greening is
  largely co-located with regions that were already moss-bearing. This supports
  the interpretation of seasonal color change rather than rapid colonization of
  entirely new zones.}
  \label{fig:persistent-and-greening}
\end{figure}

Table~\ref{tab:summary} summarizes the main spatial statistics. The median
lateral position of stable moss is \(\SI{3.7}{m}\), placing the stable moss core
well away from the centerline. The corresponding median trough depth is only
\(\SI{0.088}{m}\), indicating that the stable zone is associated with shallow to
intermediate topographic depression rather than the deepest drainage core.

\begin{table}[H]
  \centering
  \caption{Summary metrics for the stable moss pattern.}
  \label{tab:summary}
  \begin{tabular}{>{\raggedright\arraybackslash}p{9.5cm}r}
    \toprule
    Metric & Value \\
    \midrule
    Cells with stable moss signal, present on at least two dates & 26.6\% \\
    Cells strong on all three dates & 7.1\% \\
    Median lateral position of the stable moss zone, \(|x|\) & \SI{3.7}{m} \\
    Median trough depth of the stable moss zone & \SI{0.088}{m} \\
    Median \(\log(1+\mathrm{flow})\) of the stable moss zone & 2.20 \\
    Median downstream position of the stable moss zone & \SI{17.65}{m} \\
    \bottomrule
  \end{tabular}
\end{table}

\section{Dendritic Moss-Free Pattern and Topographic Controls}

The rectified photographs of the LEO1 surface show that regions of weak or
absent moss are not randomly distributed. Instead, they form a connected
dendritic pattern that resembles a surface drainage network. This pattern is
consistent with preferential flow over an inclined heterogeneous substrate:
small topographic depressions collect runoff, and neighboring depressions
become connected into larger rills and tributary-like channels.

The central dark corridor and its lateral branches therefore appear to be
negative biological imprints of the most hydrologically active parts of the
surface. Although these corridors are likely wetter than the surrounding
shoulder zones, they need not provide the most favorable habitat for moss.
Concentrated flow can suppress persistent moss cover through several coupled
mechanisms:
\begin{enumerate}[leftmargin=1.5em]
	\item \textbf{Excess saturation.} Persistent water saturation can reduce
	oxygen availability in the near-surface layer and create unfavorable
	conditions for moss growth.
	
	\item \textbf{Surface flushing.} Concentrated runoff can remove spores,
	fine particles, organic material, and weakly attached biological crust.
	
	\item \textbf{Sediment redistribution.} Flowing water can repeatedly
	transport and redeposit fine sediment along active drainage paths, thereby
	disturbing the colonized surface.
	
	\item \textbf{Mechanical instability.} Active flow corridors may undergo
	more frequent surface reorganization than adjacent shoulders, reducing the
	likelihood of stable moss establishment.
	
	\item \textbf{Moisture-window limitation.} Moss growth is favored not by
	maximal wetness, but by an intermediate moisture regime. The flanks of the
	drainage network can remain sufficiently moist while avoiding prolonged
	saturation and strong flushing.
\end{enumerate}

Thus the dendritic moss-free pattern is interpreted as a drainage-controlled
exclusion zone rather than as a primary biological branching structure. In this
interpretation, moss-rich regions occupy the shoulders and flanks of the
drainage network, whereas moss-poor regions mark the most active flow
corridors. The photographs therefore support a bounded-moisture hypothesis:
persistent moss cover is favored where water availability is sufficient but not
extreme, while the deepest and most actively draining portions of the central
trough are too wet, too unstable, or too frequently flushed to support sustained
colonization.

\section{Topographic Response Structure}

The response curves in Fig.~\ref{fig:response-curves} provide quantitative
support for this interpretation. The persistent moss proxy increases with
lateral distance from the centerline, indicating that moss cover is displaced
away from the central trough. Its dependence on trough depth is non-monotone:
moss is favored at shallow to intermediate depths and decreases sharply once the
trough becomes sufficiently deep. In addition, large values of the
flow-accumulation proxy are associated with weaker moss signal.

\begin{figure}[H]
	\centering
	\safeincludegraphics[width=0.95\linewidth]{Images/moss_response_curves.jpg}
	\caption{Response curves for persistent moss and seasonal greening. The
		persistent moss proxy increases with distance from the centerline, reaches its
		largest values at shallow to intermediate trough depths, and decreases in the
		deepest trough. The decline beyond approximately \(\SI{0.4}{m}\) provides
		evidence against a monotone ``wetter is better'' interpretation.}
	\label{fig:response-curves}
\end{figure}

The same spatial niche is visible in bath coordinates
(Fig.~\ref{fig:zone-heatmap}). The highest mean persistent moss signal occurs
in the outer half of the bath, especially for
\(|x| \gtrsim \SI{3}{m}\), and is most pronounced in the middle and lower
portions of the domain.

\begin{figure}[H]
	\centering
	\safeincludegraphics[width=0.62\linewidth]{Images/moss_zone_heatmap.jpg}
	\caption{Mean persistent moss proxy as a function of lateral position
		\(|x|\) and downstream coordinate \(s\). The highest values occur away from
		the centerline and are concentrated in the middle and lower portions of the
		bath.}
	\label{fig:zone-heatmap}
\end{figure}

\section{Cross-Sectional Evidence}

Cross-sectional profiles provide the clearest geometric test of the proposed
interpretation. In several sections, the elevation minimum of the central trough
is moss-poor, whereas higher moss values occur on the flanks and outer
shoulders. This demonstrates that the observed relationship is not merely an
artifact of two-dimensional map overlay or color thresholding.

\begin{figure}[H]
	\centering
	\safeincludegraphics[width=0.95\linewidth]{Images/moss_cross_sections.jpg}
	\caption{Cross-sections through the bath surface colored by the persistent
		moss proxy. The trough minima are repeatedly moss-poor, whereas stronger moss
		signal occurs on the flanks and outer shoulders. This directly supports the
		shoulder-zone interpretation.}
	\label{fig:moss-cross-sections}
\end{figure}

Together, the photographic, response-curve, and cross-sectional evidence
indicate that the stable moss habitat is not located in the deepest part of the
central drainage corridor. Instead, it is displaced toward off-channel shoulder
zones where moisture is available but hydrological disturbance is weaker.

\section{Smooth Surface Representation}

For geometric analysis, the soil surface was represented by a cubic bivariate
smoothing spline,
\[
  z(x,s) \approx S(x,s),
\]
fitted to the valid cells of \texttt{LEO1\_soil.asc}. The reported root-mean-square
error of the smoothed surface is approximately
\[
  \mathrm{RMSE} \approx \SI{0.071}{m}.
\]
The spline is used as a differentiable representation of the raster topography,
not as a physical seepage model. Its purpose is to make quantities such as
local slope, trough depth, and cross-sectional profile easier to interpret.

\section{Stable Moss Ranking Model}

Stable moss occurrence was summarized with a compact logistic ranking model,
\begin{equation}
  \logit P(M=1)
  =
  \beta_0
  + \beta_w W
  - \beta_f F
  - \gamma (d-d_0)^2
  + \beta_t T,
  \label{eq:ranking-model}
\end{equation}
where \(W\) is wall proximity, \(F\) is standardized \(\log(1+\mathrm{flow})\),
\(d\) is trough depth, and \(T\) is standardized soil thickness relative to the
container. The quadratic depth term represents a bounded moisture window: the
probability of stable moss is allowed to peak at an intermediate depth and to
decline when the trough becomes too shallow or too deep.

\begin{table}[H]
  \centering
  \caption{Fitted parameters of the stable moss ranking model.}
  \label{tab:model-params}
  \begin{tabular}{>{\raggedright\arraybackslash}p{7.0cm}r}
    \toprule
    Parameter & Estimate \\
    \midrule
    \(\beta_0\) & \(-0.307\) \\
    \(\beta_w\) & \(1.286\) \\
    \(\beta_f\) & \(0.185\) \\
    \(\gamma\) & \(11.971\) \\
    \(d_0\) & \SI{0.407}{m} \\
    \(\beta_t\) & \(0.146\) \\
    \bottomrule
  \end{tabular}
\end{table}

The model is not intended as a complete mechanistic law. It is a low-dimensional
ranking model that tests whether topographic variables can order favorable and
unfavorable regions. Its performance is moderate but meaningful: the reported
training AUC is 0.724 and the spatial cross-validated AUC is 0.657.

\begin{figure}[H]
  \centering
  \begin{subfigure}[t]{0.31\linewidth}
    \centering
    \safeincludegraphics[width=\linewidth]{Images/observed_stable_moss_zone.jpg}
    \caption{Observed stable moss}
  \end{subfigure}
  \hfill
  \begin{subfigure}[t]{0.31\linewidth}
    \centering
    \safeincludegraphics[width=\linewidth]{Images/predicted_stable_moss_probability.jpg}
    \caption{Predicted probability}
  \end{subfigure}
  \hfill
  \begin{subfigure}[t]{0.31\linewidth}
    \centering
    \safeincludegraphics[width=\linewidth]{Images/stable_moss_residual.jpg}
    \caption{Observed minus predicted}
  \end{subfigure}
  \caption{Spatial diagnostics for the stable moss ranking model. The model
  reproduces the broad displacement of stable moss away from the central hollow,
  although residual structure indicates that additional biological or
  hydrological variables remain unresolved.}
  \label{fig:spatial-model-diagnostics}
\end{figure}

\begin{figure}[H]
  \centering
  \safeincludegraphics[width=0.70\linewidth]{Images/stable_moss_calibration.jpg}
  \caption{Calibration of the ranking model. Across predicted-suitability
  deciles, the observed stable-moss fraction increases from about 7\% to about
  48\%. The monotone ordering supports the use of the model as a suitability
  ranking, even though it is not a full process model.}
  \label{fig:calibration}
\end{figure}

\begin{figure}[H]
  \centering
  \begin{subfigure}[t]{0.47\linewidth}
    \centering
    \safeincludegraphics[width=\linewidth]{Images/phase_curve_trough_depth_optimum.jpg}
    \caption{Trough-depth response}
  \end{subfigure}
  \hfill
  \begin{subfigure}[t]{0.47\linewidth}
    \centering
    \safeincludegraphics[width=\linewidth]{Images/phase_curve_wall_effect.jpg}
    \caption{Wall-proximity response}
  \end{subfigure}
  \caption{One-dimensional model responses. The depth curves exhibit a
  non-monotone optimum near \(\SI{0.4}{m}\), while wall proximity increases the
  predicted probability over the range shown. Increasing flow shifts suitability
  downward.}
  \label{fig:model-response-curves}
\end{figure}

\section{Biophysical Interpretation}

The observed pattern is most naturally interpreted as evidence for a bounded
moisture optimum. The central trough is likely wetter than the flanks, but it
may also experience longer saturation, lower oxygen availability, sediment
redistribution, flushing, and more frequent surface disturbance. These effects
can make the wettest zone less favorable biologically, even if water is not
limiting there.

The flanks of the trough provide a plausible intermediate regime. They remain
moist but are less likely to experience persistent waterlogging or strong
surface washout. The positive wall-proximity term in the ranking model suggests
that the box-like container geometry may further stabilize these outer shoulder
regions, either by retaining moisture laterally or by reducing erosive surface
flow.

This interpretation is consistent with broader observations from moss and
biological soil-crust systems: moisture promotes activity only up to an optimum,
whereas excessive saturation, hypoxic stress, runoff, and disturbance can
suppress growth or persistence \cite{Liu2026,Bartels2018,Macek2022,Szyja2023,Review2023}.

\section{Limitations and Needed Measurements}

The present analysis establishes a spatial association between moss persistence
and bath geometry. It does not directly measure the water state, oxygen status,
or disturbance history of the substrate. A mechanistic test of the bounded
moisture hypothesis would require at least one of the following:
\begin{itemize}[leftmargin=1.5em]
  \item local time series of soil moisture or matric potential;
  \item observations of standing water or saturation duration in the trough;
  \item tracer experiments to quantify flushing and residence time;
  \item repeated high-resolution surveys after controlled irrigation events;
  \item measurements of fine-sediment redistribution and surface stability.
\end{itemize}

These measurements would distinguish between two related mechanisms: excessive
wetness itself and hydrologic instability caused by concentrated flow.

\section{Conclusion}

The dominant biological pattern in the LEO1 bath is not simple occupation of
the wettest available ground. The stable moss footprint is displaced away from
the deepest central trough and is concentrated on off-center shoulder zones.
The combined evidence from rectified imagery, response curves, cross-sections,
and a compact ranking model supports a bounded-moisture interpretation: moss is
favored where the substrate remains sufficiently moist but is not maximally wet,
strongly flushed, or hydrologically unstable.

\appendix
\section{Additional Phase-Model Surfaces}

The following panels show the model probability surfaces under different flow
conditions. They are included as supporting diagnostics; the main text relies
primarily on the observed footprint, cross-sections, response curves, and
one-dimensional model responses.

\begin{figure}[H]
  \centering
  \begin{subfigure}[t]{0.32\linewidth}
    \centering
    \safeincludegraphics[width=\linewidth]{Images/phase_surface_simple_low_flow.jpg}
    \caption{Low flow}
  \end{subfigure}
  \hfill
  \begin{subfigure}[t]{0.32\linewidth}
    \centering
    \safeincludegraphics[width=\linewidth]{Images/phase_surface_simple_median_flow.jpg}
    \caption{Median flow}
  \end{subfigure}
  \hfill
  \begin{subfigure}[t]{0.32\linewidth}
    \centering
    \safeincludegraphics[width=\linewidth]{Images/phase_surface_simple_high_flow.jpg}
    \caption{High flow}
  \end{subfigure}
  \caption{Simple phase-model suitability surfaces under increasing flow.}
  \label{fig:simple-phase-surfaces}
\end{figure}

\begin{figure}[H]
  \centering
  \begin{subfigure}[t]{0.32\linewidth}
    \centering
    \safeincludegraphics[width=\linewidth]{Images/phase_surface_with_thickness_low_flow.jpg}
    \caption{Low flow}
  \end{subfigure}
  \hfill
  \begin{subfigure}[t]{0.32\linewidth}
    \centering
    \safeincludegraphics[width=\linewidth]{Images/phase_surface_with_thickness_median_flow.jpg}
    \caption{Median flow}
  \end{subfigure}
  \hfill
  \begin{subfigure}[t]{0.32\linewidth}
    \centering
    \safeincludegraphics[width=\linewidth]{Images/phase_surface_with_thickness_high_flow.jpg}
    \caption{High flow}
  \end{subfigure}
  \caption{Phase-model suitability surfaces with the soil-thickness term. The
  additional term modestly changes the ranking but does not alter the main
  interpretation: stable moss is favored away from the deepest central trough.}
  \label{fig:thickness-phase-surfaces}
\end{figure}

\begin{thebibliography}{9}

\bibitem{Liu2026}
Liu, L., Huang, B., Lai, G., Zhou, W., Wu, N., Xiao, Q., Wu, L., Xiong, J., et al. (2026).
Water table-temperature synergy drives maximal growth and photosynthesis in subtropical
\textit{Sphagnum palustre} under moderate regimes.
\textit{BMC Plant Biology}, 26, 318.
\url{https://url.usb.m.mimecastprotect.com/s/DU0gCLAmQOfX4LXmmSBf2TyLYhR?domain=pmc.ncbi.nlm.nih.gov

\bibitem{Bartels2018}
Bartels, S.~F., et al. (2018).
Bryophyte abundance and community composition in humid forests:
relations with depth-to-water and moisture regime.
\textit{Frontiers in Plant Science}, 9, 858.
\url{https://url.usb.m.mimecastprotect.com/s/SE6iCM7nRPI9Wj9RRskhXT8HgjI?domain=frontiersin.org

\bibitem{Macek2022}
Macek, M., et al. (2022).
Topography cannot fully substitute microclimate for bryophyte ecology.
\textit{Science of the Total Environment}, 822, 153377.
\url{https://url.usb.m.mimecastprotect.com/s/zR2HCN7oVgI9Xo9EEsjilTyvmgF?domain=doi.org

\bibitem{Szyja2023}
Szyja, M., Felde, V.~J.~M.~N.~L., L\"uckel, S., Tabarelli, M., Leal, I.~R., B\"udel, B., Wirth, R. (2023).
Biological soil crusts decrease infiltration but increase erosion resistance in a human-disturbed tropical dry forest.
\textit{Frontiers in Microbiology}, 14, 1136322.
\url{https://url.usb.m.mimecastprotect.com/s/1hVFCOJpWjfw43w00IrsKTGej1j?domain=pubmed.ncbi.nlm.nih.gov

\bibitem{Review2023}
Mata-Gonz\'alez, R., et al. (2023).
A review on effects of biological soil crusts on hydrological processes.
\textit{Earth-Science Reviews}, 243, 104516.
\url{https://url.usb.m.mimecastprotect.com/s/LO4KCP6q0kCZMQZooh6tmTxijgn?domain=sciencedirect.com

\end{thebibliography}

\end{document}

